\begin{small}
\begin{longtable}{p{1.0cm}p{7.6cm}rr}
\caption{Туршилтын үр дүн (dev, кирилл асуулт, EX \%)}\label{tab:results-dev}\\
\toprule ID & Тохиргоо & Arctic-R1-7B & sqrl-9b \\\midrule\endfirsthead
\toprule ID & Тохиргоо & Arctic-R1-7B & sqrl-9b \\\midrule\endhead
\multicolumn{4}{l}{\textbf{A. Суурь загвар, шууд (сургалтгүй, RAG-гүй)}}\\
A1 & Кирилл асуулт, тайлбаргүй түүхий схем & — & — \\
A2 & A1-ийн англи хувилбар (хэлний зөрүү = A2 − A1) & — & — \\
A3 & A1-ийн латин галиг хувилбар (галигийн зөрүү = A1 − A3) & — & — \\
A4 & Асуултыг англи руу орчуулах, утгыг placeholder-оор хамгаалах & — & — \\
A5 & DB судлах алхамгүй (max\_steps = 0) & n/a & — \\
A6 & Зорилтот dialect-ийг (PostgreSQL) prompt-д заах & — & — \\
A7 & SQLite-аар үүсгээд sqlglot-оор PostgreSQL руу хөрвүүлэх & — & — \\
\multicolumn{4}{l}{\textbf{B. Суурь загвар + RAG ба few-shot}}\\
B1 & Хүснэгт, баганын монгол/англи тайлбар, синоним (+ өнөөдөр, валютын дүрэм) & — & — \\
B2 & Утгын хайлт (кирилл, галиг түлхүүр), нийт 10–20 утга & — & — \\
B3 & Бизнес нэр томьёо (evidence) & — & — \\
B4 & 3–5 few-shot жишээ (DAIL маягаар) & — & — \\
B5 & B1–B4 бүгд (RAG-full) & — & — \\
B6 & B5 − тайлбар (B1-ийн бүх хэсэг) & — & — \\
B7 & B5 − утгын хайлт & — & — \\
B8 & B5 − нэр томьёо & — & — \\
B9 & B5 − few-shot & — & — \\
B10 & B5 + schema linking (top-k хүснэгт) & — & — \\
\multicolumn{4}{l}{\textbf{C. Зөвхөн fine-tune}}\\
C1 & Rejection sampling SFT (8 гаргалт, гүйцэтгэл таарсан) & — & — \\
C2 & C1 + gold SQL-ээр чиглүүлсэн CoT & — & — \\
C3 & C2 + өөрийн DB-ийн синтетик ба жишээний сан (160) & — & — \\
C4 & C3-ийн өгөгдөл, англи 30\% (C4b: 50\%) & — & — \\
C4b & C3-ийн өгөгдөл, англи 50\% & — & — \\
C5 & C3, LoRA rank 64 (C3 = rank 32) & — & — \\
C6 & C-ийн шилдэг SFT + GRPO (1 / 0.1 / 0, 8 rollout) & — & — \\
C7 & C6 + англи асуултын replay & — & — \\
\multicolumn{4}{l}{\textbf{D. Fine-tune + RAG}}\\
D1 & C-ийн шилдэг загвар + B5-ийн RAG, few-shot & — & — \\
D2 & B5-ийн RAG-тай prompt-оор сургаж RAG-тай ашиглах & — & — \\
D3 & D2 − few-shot & — & — \\
D4 & D2 + GRPO (RAG-тай prompt) & — & — \\
\multicolumn{4}{l}{\textbf{E. Засвар, санал хураалт, татгалзалт}}\\
E1 & Алдааны мессежээр засах (2 удаа хүртэл; төлөвлөгөө: 1–3) & — & — \\
E2 & 8 нэр дэвшигч SQL-ийн санал хураалт (гүйцэтгэлийн үр дүнгээр) & — & — \\
E3 & Санал нийлэлтээр татгалзах (босгыг E2-ийн муруйгаас сонгоно) & — & — \\
\multicolumn{4}{l}{\textbf{F. Монгол хэлэнд чиглэсэн санаа}}\\
F1 & Монгол BIRD хувилбар дээр A1 & — & — \\
F1b & Монгол BIRD хувилбар дээр B5 (BIRD evidence) & — & — \\
F2 & B5 + дүрмийн бичиг хоорондын утга тааруулалт (дагавар, галиг) & — & — \\
F3 & F2 + хэрэглэгчийн логоос сургасан галиг хувиргагч & — & — \\
F4 & C6 + бичгээс үл хамаарах reward & — & — \\
F5 & A + англи pivot дахин томьёолол (prompt-оор) & — & — \\
F6 & C-ийн шилдэг жор + англи pivot-той сургалтын мөр & — & — \\
F7 & Монгол хэлэнд хүчтэй жижиг загвар → intent JSON → SQL (A4-тэй харьцуулна) & — & — \\
F8 & Semantic layer + чиглүүлэгч (KPI-г шууд compile) & — & — \\
F9a & Санал зөрөлтөөр сонгосон 100 тэмдэглэгээг жишээний санд нэмсэн & — & — \\
F9b & Санамсаргүй сонгосон 100 тэмдэглэгээг нэмсэн (F9a-тай харьцуулна) & — & — \\
F10 & sqrl + search\_values tool (tool-той траектороор сургасан, D-ийн жор) & n/a & — \\
\multicolumn{4}{l}{\textbf{G. Agentic систем (төлөвлөгөөнөөс гадуурх нэмэлт; docs/AGENT.md)}}\\
G1 & Agent: гүйцэтгэл ба литерал шалгах tool → LLM critic → засах (≤2 удаа) & — & — \\
G2 & G1 − LLM critic (зөвхөн детерминист tool-ийн олдвороор засна) & — & — \\
G3 & G1 + planner (SQL-ээс өмнө англи төлөвлөгөө) & — & — \\
G4 & RAG-гүй agent: A-шилдэг + G1-ийн давталт & — & — \\
G5 & G1 − литерал tool (утга шалгах DB tool-ийн нэмэр = G1 − G5) & — & — \\
G6 & G2 − литерал tool = зөвхөн гүйцэтгэлийн алдаагаар нэмэлт засвар (agent-гүй давталтын суурь) & — & — \\
\bottomrule
\multicolumn{4}{p{13.4cm}}{\footnotesize F1, F1b нь монгол BIRD dev дээр (ББСБ-ийн eval set биш).}\\
\end{longtable}
\end{small}
